% =========================================================
%  sesi/07/sesi.tex
%  Sesi 7 — Mesin Keadaan Penjaga
% =========================================================

\sesi{07}{Mesin Keadaan Penjaga}{2.2}{140 menit}

Pada sesi 6 penjaga hanya mengenal dua kemungkinan: mengejar atau berdiam.
Keduanya disimpan pada satu variabel bernilai benar atau salah. Pada minggu 7
variabel itu diganti susunan yang lebih teratur. Penjaga memperoleh tiga
keadaan yang berbeda---berdiam, berpatroli, dan mengejar---beserta aturan
perpindahan di antaranya. Ia berjalan mengelilingi lorong, mengejar ketika
pemain terlihat, dan menyerah setelah kehilangan jejak selama waktu tertentu.

Perubahan itu bukan sekadar menambah cabang pemeriksaan. Ada susunan yang
memberi nama pada setiap perilaku, memisahkan syarat perpindahan dari tindakan,
dan menyimpan waktu. Susunan itulah yang disebut \istilah{Finite State Machine}
(FSM), dan padanya Anda belajar bahwa perilaku yang tampak rumit dapat diuraikan
menjadi beberapa keadaan sederhana yang saling berpindah menurut aturan yang
dapat diperiksa.

Sesi ini menuntaskan Sub-CPMK 2.2: mengimplementasikan FSM untuk NPC, dengan
transisi berbasis jarak dan \istilah{cooldown}.

% ---------------------------------------------------------
\bagian{Tujuan Sesi}

\begin{tujuanbox}
Setelah menuntaskan sesi ini, Anda dapat menunjukkan:
\begin{enumerate}[leftmargin=1.5em, itemsep=0.2em]
  \item penjaga yang berubah warna menurut keadaannya---abu-abu, hijau, dan
        merah---dengan keterangan keadaan pada HUD;
  \item diagram keadaan--transisi tiga keadaan yang setiap anak panahnya diberi
        syarat yang dapat diperiksa program;
  \item penjelasan perbedaan \istilah{state}, \istilah{event}, dan
        \istilah{transition} beserta peran masing-masing pada penjaga;
  \item penjaga yang berdiam 1,5 detik, berpatroli pada kecepatan 90 piksel per
        detik, mengejar pada 165 piksel per detik, dan menyerah setelah
        1,0 detik kehilangan jejak;
  \item penjelasan mengapa keputusan keadaan dipisahkan dari pelaksanaan
        geraknya, beserta keuntungan pemisahan itu.
\end{enumerate}
\end{tujuanbox}

% ---------------------------------------------------------
\bagian{Prasyarat}

\begin{prasyaratbox}
\begin{itemize}[leftmargin=1.4em, itemsep=0.25em]
  \item Checkpoint sesi 6 berjalan: level Lorong tampil, pemain bergerak dengan
        WASD, dan penjaga mengejar berdasarkan jarak.
  \item Anda memahami perhitungan $\vec{u} = \vec{p} - \vec{n}$ dan
        penormalan vektor yang dipakai penjaga.
  \item Anda dapat membaca kelas Python dengan metode dan atribut, serta
        memakai \code{delta_seconds} untuk menjalankan pewaktu.
  \item Bab IX buku ajar sudah dibaca, khususnya bagian tentang diagram
        keadaan dan tiga jenis pewaktu.
\end{itemize}
\end{prasyaratbox}

% ---------------------------------------------------------
\bagian{Peta Kemampuan}

\petasesi
  {Penjaga dengan satu variabel \code{chasing}; hanya mengejar atau berdiam.}
  {Tiga keadaan dengan aturan perpindahan, dua pewaktu, jalur patroli, warna
   keadaan, dan HUD keadaan.}
  {Pencarian jalur A* pada sesi 8, lalu penggabungan FSM dan pencarian jalur
   pada sesi 9.}

% ---------------------------------------------------------
\bagian{Konsep yang Dipakai}

\subsection*{Tiga kata yang harus dibedakan}

Sebuah \istilah{state} adalah cara penjaga berperilaku untuk sementara waktu:
berdiam, berpatroli, atau mengejar. Selama berada pada satu keadaan, penjaga
mengikuti aturan keadaan itu saja. Sebuah \istilah{event} adalah keterangan
yang dibaca program, misalnya jarak terhadap pemain atau habisnya sebuah
pewaktu. Sebuah \istilah{transition} adalah aturan yang memindahkan penjaga
dari satu keadaan ke keadaan lain ketika event tertentu terpenuhi. FSM adalah
kumpulan keadaan, kumpulan transisi, dan pekerjaan penyiapan yang dijalankan
saat sebuah keadaan dimasuki.

\subsection*{Tiga keadaan penjaga}

Keadaan \istilah{idle} berarti berdiri diam sambil menghitung waktu. Setelah
1,5 detik, penjaga mulai berpatroli. Keadaan \istilah{patrol} berarti berjalan
dari satu titik patroli ke titik berikutnya pada kecepatan 90 piksel per detik.
Keadaan \istilah{chase} berarti bergerak menuju pemain pada kecepatan 165
piksel per detik, dengan arah yang dihitung seperti pada sesi 6.

\begin{center}% figures/diagram-keadaan.tikz
% Tiga keadaan penjaga dan empat perpindahan yang dipakai checkpoint.
\begin{tikzpicture}[
  keadaan/.style={rectangle, rounded corners=8pt, draw=ubNavy, thick,
    fill=ubSoft, align=center, text width=2.2cm, minimum height=0.95cm,
    font=\small\bfseries},
  serang/.style={keadaan, draw=ubRed, fill=ubRedBg},
  panah/.style={-{Stealth[length=2.4mm]}, thick, draw=ubAccent},
  label/.style={font=\scriptsize, fill=white, inner sep=1.5pt}
]
  \node[keadaan] (idle) at (0,0) {berdiam};
  \node[keadaan] (patrol) at (5.0,0) {berpatroli};
  \node[serang] (chase) at (2.5,-2.5) {mengejar};

  \draw[panah] (idle) to[bend left=12] node[label, above] {pewaktu habis} (patrol);

  \draw[panah] (patrol) to[bend left=10]
    node[label, right=2pt, midway] {jarak $\le$ radius} (chase);
  \draw[panah] (chase) to[bend left=10]
    node[label, left=2pt, midway] {hilang $>$ 1 detik} (patrol);

  \draw[panah, dashed] (idle) to[bend right=10]
    node[label, left=2pt, midway] {pemain terlihat} (chase);
\end{tikzpicture}
\end{center}

Diagram di atas memperlihatkan tiga keadaan dan empat perpindahan yang dipakai
checkpoint ini: berdiam menuju berpatroli ketika pewaktu berdiam habis,
berpatroli menuju mengejar ketika jarak pemain tidak lebih dari 252 piksel,
mengejar menuju berpatroli setelah pemain tidak terlihat selama 1,0 detik, dan
berdiam menuju mengejar jika pemain terlihat lebih dahulu. Perpindahan terakhir
mencegah penjaga menunggu pewaktu berdiam habis ketika pemain muncul dekat
dengannya.

\subsection*{Pewaktu sebagai syarat perpindahan}

Dua perpindahan pada sesi ini tidak ditentukan oleh jarak, melainkan oleh waktu.
Pewaktu berdiam mulai dari 1,5 detik dan berkurang setiap bingkai; ketika
nilainya mencapai nol, penjaga berpatroli. Pewaktu menyerah mulai dari nol dan
bertambah selama pemain berada di luar radius; ketika mencapai 1,0 detik,
penjaga kembali berpatroli. Selama pemain berada di dalam radius, pewaktu itu
dikembalikan ke nol, sehingga penjaga tidak menyerah hanya karena pemain
berlari di balik dinding selama satu bingkai.

Kata \istilah{cooldown} dipakai untuk tenggang waktu yang harus dilalui sebelum
sebuah keputusan boleh berubah. Pada sesi ini perannya dipegang kedua pewaktu
di atas. Perhatikan bahwa urutan pemeriksaan penting: pada keadaan berdiam,
syarat jarak diperiksa lebih dahulu daripada syarat pewaktu, agar penjaga yang
melihat pemain segera mengejar walaupun baru saja mulai berdiam.

\subsection*{Keputusan dipisahkan dari tindakan}

Seluruh syarat perpindahan dikumpulkan pada kelas \code{GuardBrain}, sedangkan
gerak badan penjaga dikerjakan kelas \code{Guard}. \code{GuardBrain} membaca
jarak dan kedua pewaktu, lalu mengembalikan nama keadaan; ia tidak memindahkan
apa pun. \code{Guard} menerima nama keadaan itu, mengganti warna, dan berjalan
menurut kecepatan keadaan tersebut. Pemisahan itu membuat aturan perpindahan
dapat diperiksa tanpa menjalankan permainan, dan membuat penambahan keadaan
baru hanya menuntut satu blok baru pada \code{GuardBrain}.

% ---------------------------------------------------------
\bagian{Tutorial}

Kerjakan langkah berikut berurutan. Simpan proyek sesi 6 sebagai salinan baru
sebelum mulai, agar checkpoint lama tetap dapat dibandingkan.

\begin{langkah}
  \item \textbf{Siapkan proyek.} Buka folder sesi 6, jalankan
        \code{python main.py --level 2}, dan pastikan penjaga masih mengejar
        pemain. Salin folder itu menjadi folder kerja sesi 7.

  \item \textbf{Tambahkan pengaturan.} Pada \file{settings.py}, tambahkan warna
        keadaan \code{GUARD_PATROL_COLOR}, lalu tambahkan
        \code{GUARD_IDLE_SECONDS = 1.5}, \code{GUARD_GIVE_UP_SECONDS = 1.0},
        dan \code{WAYPOINT_REACHED = 6.0}. Hapus
        \code{GUARD_STOP_DISTANCE} apabila sudah tidak dipakai.

  \item \textbf{Tentukan titik patroli.} Pada \file{levels.py}, buat konstanta
        \code{PATROL_POINTS} berisi empat petak, misalnya
        $(4,2)$, $(10,2)$, $(10,12)$, dan $(4,12)$. Tambahkan kunci
        \code{waypoints} pada kamus yang dikembalikan \code{build_lorong}.

  \item \textbf{Buat berkas mesin keadaan.} Buat \file{gudang_artefak/fsm.py}.
        Isinya tiga konstanta nama keadaan, kelas \code{GuardBrain}, dan metode
        \code{update}, \code{patrol_point}, serta \code{next_patrol_point}.

  \item \textbf{Tulis keputusan keadaan.} Pada \code{update}, hitung jarak ke
        pemain. Periksa keadaan yang sedang aktif satu per satu: dari berdiam
        ke mengejar bila jarak memenuhi radius, dari berdiam ke berpatroli bila
        pewaktu habis, dan dari mengejar ke berpatroli bila pewaktu menyerah
        habis. Kembalikan nama keadaan pada akhir fungsi.

  \item \textbf{Tulis gerak patroli.} \code{patrol_point} mengubah petak tujuan
        menjadi piksel dengan rumus titik pusat petak. \code{next_patrol_point}
        memajukan indeks dan kembali ke titik pertama setelah titik terakhir.

  \item \textbf{Ubah kelas penjaga.} Ganti atribut \code{chasing} dengan
        atribut \code{state}. Tambahkan \code{set_state} yang mengganti warna
        badan menurut keadaan, dan \code{walk_to} yang berjalan lurus menuju
        titik serta mengembalikan benar bila sudah tiba.

  \item \textbf{Sambungkan ke permainan.} Pada \file{game.py}, buat
        \code{GuardBrain} ketika level Lorong dimuat. Setiap bingkai, minta
        keputusan, panggil \code{set_state}, lalu jalankan keadaan itu: mengejar
        memakai \code{walk_to} ke pemain, berpatroli memakai
        \code{walk_to} ke titik patroli. Ketika pemain tertangkap, buat ulang
        otak penjaga agar kembali ke keadaan awal.

  \item \textbf{Tambahkan HUD keadaan.} Pada \code{_draw_hud}, tampilkan
        keadaan penjaga di sudut kanan atas, misalnya
        \code{Penjaga: chase}. Warna teks dibuat netral agar tidak tertukar
        dengan warna badan penjaga.

  \item \textbf{Jalankan dan ukur.} Jalankan
        \code{python main.py --level 2}. Amati warna penjaga, keterangan HUD,
        dan jalur patrolinya. Ukur dengan pengukur waktu: lama berdiam, lama
        mengejar sebelum menyerah, dan lama satu putaran patroli.
\end{langkah}

\begin{ujibox}
Setelah Langkah 10, periksa lima hal berikut.
\begin{itemize}[leftmargin=1.4em, itemsep=0.2em]
  \item Penjaga berwarna abu-abu ketika program dimulai, lalu berubah hijau
        setelah sekitar 1,5 detik. Bila ia langsung hijau, periksa nilai awal
        \code{idle_timer}.
  \item Penjaga berubah merah ketika pemain berada dalam 252 piksel, dan HUD
        menampilkan keadaan yang sama dengan warnanya.
  \item Penjaga kembali hijau sekitar 1,0 detik setelah pemain keluar dari
        radius. Bila ia menyerah seketika, periksa pengembalian
        \code{lost_timer} ke nol.
  \item Penjaga berjalan melalui keempat titik patroli dan mengulanginya dari
        titik pertama, dengan kecepatan yang terasa lebih lambat daripada
        kecepatan mengejar.
  \item Garis merah hanya tampil pada keadaan mengejar. Bila garis tampil pada
        keadaan lain, periksa syarat penggambarannya.
\end{itemize}
\end{ujibox}

% ---------------------------------------------------------
\bagian{Kode Sesi Ini}

Potongan berikut adalah bagian pokok checkpoint sesi 7. Berkas lengkapnya
tersedia pada folder \file{code/sesi-07}. Perhatikan bahwa kelas penjaga pada
sesi ini tidak lagi menghitung arah sendiri; pekerjaan itu berpindah ke
\code{GuardBrain} dan ke metode \code{walk_to}.

\tampilkodebaris[Python]{code/sesi-07/gudang_artefak/settings.py}{49}{56}
  {Pengaturan penjaga dan mesin keadaan}

\tampilkodebaris[Python]{code/sesi-07/gudang_artefak/levels.py}{72}{76}
  {Titik patroli dan titik awal penjaga}

\tampilkodebaris[Python]{code/sesi-07/gudang_artefak/fsm.py}{14}{58}
  {Mesin keadaan: keadaan, transisi, dan titik patroli}

\tampilkodebaris[Python]{code/sesi-07/gudang_artefak/entities.py}{208}{228}
  {Mengubah warna menurut keadaan dan berjalan menuju titik}

\tampilkodebaris[Python]{code/sesi-07/gudang_artefak/game.py}{141}{177}
  {Meminta keputusan, menjalankan keadaan, dan patroli}

\tampilkodebaris[Python]{code/sesi-07/gudang_artefak/game.py}{236}{245}
  {HUD yang menampilkan keadaan penjaga}

\checkpoint{07}{Program pada folder itu dapat dijalankan dengan
\code{python main.py --level 2}; penjaga berdiam 1,5 detik, berpatroli pada
kecepatan 90 piksel per detik, mengejar pada 165 piksel per detik, dan menyerah
setelah 1,0 detik kehilangan jejak.}

% ---------------------------------------------------------
\bagian{Kesalahan Umum}

\begin{galatbox}
\begin{itemize}[leftmargin=1.4em, itemsep=0.35em]
  \item \textbf{Penjaga tidak pernah meninggalkan keadaan berdiam.}
        Gejala: warna tetap abu-abu dan pewaktu tidak berjalan. Penyebabnya,
        \code{idle_timer} dimulai dari nol atau tidak dikurangi. Periksa nilai
        awal \code{GUARD_IDLE_SECONDS} dan baris
        \code{self.idle_timer -= delta_seconds}.
  \item \textbf{Penjaga berkedip antara mengejar dan berpatroli.}
        Gejala: warna berganti-ganti ketika pemain berada tepat di tepi radius.
        Penyebabnya, pewaktu menyerah tidak dikembalikan ke nol selama pemain
        masih terlihat. Periksa cabang \code{if jarak <= self.radius}.
  \item \textbf{Penjaga menyerah seketika.}
        Gejala: penjaga berhenti mengejar pada bingkai pertama setelah pemain
        menjauh. Penyebabnya, \code{lost_timer} tidak direset ketika keadaan
        mengejar dimasuki. Periksa \code{self.lost_timer = 0.0}.
  \item \textbf{Warna tidak berubah.}
        Gejala: HUD berubah tetapi badan penjaga tetap abu-abu. Penyebabnya,
        \code{set_state} keluar lebih awal karena nama keadaan dianggap sama,
        atau kunci kamus warna tidak lengkap. Periksa ketiga kunci
        \code{idle}, \code{patrol}, dan \code{chase}.
  \item \textbf{Penjaga berhenti di antara dua titik patroli.}
        Gejala: penjaga diam di tempat dan tidak melanjutkan perjalanan.
        Penyebabnya, ambang \code{WAYPOINT_REACHED} lebih kecil daripada
        perpindahan satu bingkai, sehingga titik tidak pernah dianggap tercapai.
        Periksa nilai 6.0 piksel dan kecepatan patroli 90 piksel per detik.
  \item \textbf{\code{KeyError: 'waypoints'}.}
        Gejala: level Lorong gagal dimuat. Penyebabnya, kamus hasil
        \code{build_lorong} belum memuat kunci \code{waypoints}. Periksa
        pengembalian fungsi itu dan nama kunci pada \code{_load_lorong}.
  \item \textbf{HUD bertumpuk dengan skor.}
        Gejala: keterangan keadaan menimpa teks skor. Penyebabnya, posisi
        horizontal dihitung dari lebar jendela tanpa menyisakan ruang untuk
        teks. Geser posisi itu ke kiri atau pendekkan keterangannya.
\end{itemize}
\end{galatbox}

% ---------------------------------------------------------
\clearpage
\bagian{Lembar Kerja}

Isi tabel berikut pada catatan praktikum Anda. Lakukan pengukuran dengan
pengukur waktu atau dengan menghitung bingkai pada layar, lalu bandingkan
hasilnya dengan nilai pada \file{settings.py}.

{\small
\renewcommand{\arraystretch}{1.4}
\begin{tabularx}{\linewidth}{@{}C{0.8cm} L{6.2cm} Y@{}}
\toprule
\textbf{No} & \textbf{Hasil yang dicatat} & \textbf{Keterangan Anda} \\
\midrule
1 & Lama penjaga berdiam sebelum berpatroli & \\[0.55cm]
2 & Lama penjaga mengejar sebelum menyerah & \\[0.55cm]
3 & Lama satu putaran patroli dan banyak titik yang dilalui & \\[0.55cm]
4 & Warna penjaga pada setiap keadaan & \\[0.55cm]
5 & Keadaan yang paling sulit Anda amati, beserta alasannya & \\[0.9cm]
6 & Satu kendala yang Anda hadapi dan cara menyelesaikannya & \\[0.9cm]
\bottomrule
\end{tabularx}
}

% ---------------------------------------------------------
\bagian{Checklist Sesi}

\begin{checklistbox}
\begin{checklist}
  \item Saya dapat menunjukkan ketiga keadaan penjaga melalui warna dan HUD.
  \item Saya dapat menggambar diagram keadaan--transisi dengan syarat yang
        dapat diperiksa pada setiap anak panah.
  \item Saya dapat menjelaskan perbedaan \istilah{state}, \istilah{event}, dan
        \istilah{transition}.
  \item Saya dapat menjelaskan peran pewaktu berdiam 1,5 detik dan pewaktu
        menyerah 1,0 detik.
  \item Saya dapat menjelaskan mengapa syarat jarak diperiksa sebelum syarat
        pewaktu pada keadaan berdiam.
  \item Saya dapat menjelaskan keuntungan memisahkan \code{GuardBrain} dari
        kelas \code{Guard}.
\end{checklist}
\end{checklistbox}

% ---------------------------------------------------------
\bagian{Pengait ke Sesi Berikutnya}

Mesin keadaan menjawab pertanyaan ``apa yang sebaiknya penjaga lakukan?''. Ia
belum menjawab pertanyaan ``bagaimana cara menuju ke sana?''. Pada keadaan
mengejar, penjaga masih berjalan lurus ke arah pemain, sehingga dinding tetap
menghalanginya. Sesi 8 menambahkan pencarian jalur A* pada peta petak yang sudah
Anda bangun. Keadaan mengejar tidak lagi memakai arah langsung, melainkan rute
yang dihitung dari petak awal menuju petak pemain. Dengan begitu, lapisan
keputusan yang Anda susun pada sesi ini akhirnya berhubungan dengan lapisan
navigasi.
